% lecture14.tex: Lecture 14: Synthesis, Solar System in Context, Astrobiology
% Planetary Systems, Kapteyn Institute

\documentclass[aspectratio=169, 11pt]{beamer}

% ── Theme and macros ────────────────────────────────────────
\usepackage{../common/beamerthemeIPS}
% macros.tex: Shared math macros for IPS lecture slides
% Mirrors definitions in book/_config.yml for consistency
% Uses \providecommand to avoid conflicts with unicode-math

% ── Derivatives ─────────────────────────────────────────────
\providecommand{\dv}[2]{\frac{\mathrm{d} #1}{\mathrm{d} #2}}
\providecommand{\dd}{}\renewcommand{\dd}{\mathrm{d}}
\providecommand{\pdv}[2]{\frac{\partial #1}{\partial #2}}

% ── Solar quantities ────────────────────────────────────────
\newcommand{\Msun}{M_{\odot}}
\newcommand{\Rsun}{R_{\odot}}
\newcommand{\Lsun}{L_{\odot}}

% ── Planetary quantities ────────────────────────────────────
\newcommand{\Mearth}{M_{\oplus}}
\newcommand{\Rearth}{R_{\oplus}}
\newcommand{\Mjup}{M_{\mathrm{J}}}
\newcommand{\Rjup}{R_{\mathrm{J}}}

% ── Constants ───────────────────────────────────────────────
\newcommand{\kB}{k_{\mathrm{B}}}

% ── Media links ─────────────────────────────────────────────
% Clickable button that opens the animated or video version of a
% figure, e.g. \mediabutton{https://...gif}{Play animation}
\providecommand{\mediabutton}[2]{\href{#1}{\beamergotobutton{#2}}}

\graphicspath{{figures/}{../common/}{../../book/14_synthesis/figures/}}

% ── Metadata ────────────────────────────────────────────────
\title[Lecture 14: Synthesis]{Lecture 14: Synthesis\\Solar System in Context, Astrobiology}
\subtitle{Planetary Systems}
\author[L14]{Tim Lichtenberg}
\institute{Kapteyn Astronomical Institute\\University of Groningen}
\date{2026}
\titlebgcredit{Tidally locked rocky planet atmospheric circulation. Wordsworth \& Kreidberg (2022)}

% ════════════════════════════════════════════════════════════
\begin{document}

% ── Title slide ─────────────────────────────────────────────
\begin{frame}[plain,noframenumbering]
  \titlepage
\end{frame}

% ── Learning objectives ─────────────────────────────────────
\begin{frame}{Learning Objectives}
  By the end of this lecture, you will be able to:
  \vskip8pt
  \begin{itemize}
    \item Place our solar system within the observed exoplanet population
    \item Evaluate habitability as a coupled \emph{systems} property and derive the classical habitable-zone boundaries
    \item Identify the most promising life-detection targets for the next decade
  \end{itemize}
  \vskip8pt
  \keyresult{Habitability is a trajectory, not a snapshot.}
\end{frame}

\begin{frame}
  \ipshero[0.76]{solar_system_composite}{The principal bodies of the solar system grouped by region at true relative size, from the terrestrial planets to the Kuiper belt objects. Credit: CactiStaccingCrane (Wikimedia Commons, CC BY-SA 4.0); NASA/ESA/ISRO.}
\end{frame}

% ── Recap ───────────────────────────────────────────────────
\begin{frame}{The Thread of the Course}
  \begin{itemize}
    \item L1--L8: accretion, differentiation, internal heat engines, and atmosphere-surface coupling
    \item L9--L12: terrestrial and giant planet evolution traced body by body
    \item L13: over 6000 exoplanets revealing architectural diversity
  \end{itemize}
  \vskip8pt
  \textbf{Today:} Step back to evaluate the solar system in context and test astrobiological habitability.
\end{frame}

\begin{frame}
  \ipshero{solar_system_architecture}{Architecture of the solar system on a logarithmic semi-major axis scale with symbol area scaled to radius; shaded spans mark the asteroid belt, the Kuiper belt and the water snow line region. Credit: Course-original figure; NASA NSSDC.}
\end{frame}


% ════════════════════════════════════════════════════════════
\sectionimage{drazkowska2023_growth_processes}{Dust-to-planet growth. Dr\k{a}\.zkowska et al.\ (2023)}
\section{Solar System in the Exoplanet Context}
% ════════════════════════════════════════════════════════════

\begin{frame}
  \ipsherokey{drazkowska2023_growth_processes}{Growth processes from dust to planets. Reproduced from Dr\k{a}\.zkowska et al.\ (2023).}{Dust grows to pebbles, planetesimals and embryos by streaming instability, pebble accretion and runaway envelope accretion; the disk lifetime of 3--5~Myr is the deadline, and ALMA disk substructure is direct observational input.}
\end{frame}

\begin{frame}
  \ipsherokey[0.61]{lambrechts2012_core_growth}{Core mass against time at 0.5, 5 and 50 AU: pebble accretion (solid) against classical planetesimal accretion (grey dotted), with Ceres and Pluto marked. Reproduced from Lambrechts \& Johansen (2012).}{Pebble accretion reaches 10 Earth masses well within a typical disk lifetime; planetesimal accretion at 5 AU takes longer than the disk lives.}
\end{frame}

\begin{frame}
  \ipsherokey[0.55]{lambrechts2012_growth_time}{Time to grow a 10 Earth mass core against distance from the star by pebble accretion (black), planetesimal accretion (grey) and fragment accretion (dashed grey); the hatched band is the 1 to 10~Myr gas dispersal. Reproduced from Lambrechts \& Johansen (2012).}{A core that is to capture an envelope must reach 10 Earth masses before the gas goes: the pebble line stays below the band at every radius, the planetesimal line rises through it within a few AU.}
\end{frame}

\begin{frame}
  \ipsherokey[0.61]{lichtenberg2023_nccc_timeline}{Timeline of solar system formation from isotopic dating of meteorites: the non-carbonaceous (NC, red) and carbonaceous (CC, blue) reservoirs. Reproduced from Lichtenberg et al.\ (2023).}{The two reservoirs keep distinct isotopic signatures from CAI formation onward, so they accreted in physically separated regions of the disk for at least the first 2 to 3~Myr.}
\end{frame}

\begin{frame}
  \ipsherokey[0.66]{andrews2018_dsharp_gallery}{ALMA DSHARP 240~GHz continuum images of 20 nearby disks. Reproduced from Andrews et al.\ (2018).}{Concentric rings, gaps and asymmetries are nearly ubiquitous, and most are interpreted as signatures of planet growth in progress.}
\end{frame}

\begin{frame}
  \ipsherokey{raymond2022_period_radius}{Transiting census with compact multi-planet systems boxed: five uniform ones (middle), three non-uniform ones (bottom). Weiss et al.\ (2023, PPVII).}{Compact multi-planet systems show striking architectural uniformity.}
\end{frame}

\begin{frame}
  \ipsherokey[0.61]{fulton2017_radius_valley}{The Fulton gap: planet size against incident stellar flux for small close-in exoplanets, with the deficit near 1.8~$\Rearth$ between super-Earths and sub-Neptunes; the lower panel overlays photoevaporation models. Reproduced from Fulton et al.\ (2017).}{Photoevaporation models reproduce both the location and the slope of the valley with irradiation.}
\end{frame}

% ── Where solar system is atypical ──────────────────────────
\begin{frame}{Where the Solar System is Atypical}
  \begin{itemize}\setlength\itemsep{4pt}
    \item No super-Earths or sub-Neptunes: the most common exoplanet size class is absent
    \item No close-in giant planets: hot Jupiters orbit $\sim 0.5$--$1\%$ of Sun-like stars, hot Neptunes are rarer still
    \item Outer giants on wide, circular orbits, while the inner terrestrials lack compact uniformity
  \end{itemize}
\end{frame}

\begin{frame}
  \ipsherokey[0.66]{drazkowska2023_population_synthesis}{Population synthesis predictions for planet mass against orbital period under planetesimal and pebble accretion. Credit: Dr\k{a}\.zkowska et al.\ (2023).}{Both pathways readily form super-Earths and close-in giants that the solar system lacks.}
\end{frame}

\begin{frame}
  \ipsherokey[0.65]{raymond2022_peas_in_a_pod}{Compact multi-systems with at least four transiting planets within 1.52~AU, ranked by the uniformity of their planet sizes. Reproduced from Weiss et al.\ (2023).}{The solar system falls in the bottom quintile of size uniformity: the inner solar system did not produce the typical compact-multi architecture.}
\end{frame}

\begin{frame}
  \ipsherokey[0.6]{lichtenberg2025_mass_radius}{Mass-radius diagram with composition tracks: pure Fe, Earth-like, MgSiO$_3$, water-rich, magma plus H$_2$O, and gas dwarf. Lichtenberg et al.\ (2025).}{The solar system terrestrials and the close-in super-Earths sit on the Earth-like curve, the latter as stripped cores that lost their gas envelopes; larger planets carry volatile envelopes that swell their radii.}
\end{frame}

\begin{frame}
  \ipsherokey[0.62]{mulders2024_occurrence_vs_teff}{Occurrence of small planets (1--4~$\Rearth$, $P < 50$~d) per star against stellar effective temperature. Mulders (2024).}{Occurrence is roughly twice as high around early M dwarfs as around F dwarfs and keeps rising into the late M dwarfs: Earth-sized planets are common around the most numerous stars.}
\end{frame}

\begin{frame}
  \ipsherokey[0.62]{elt_milkyway}{The Milky Way over the ELT site. Credit: C. Letelier/ESO.}{Whether our architecture is rare stays open: 1~AU terrestrial analogues sit at the survey detection limits, Kepler's four-year baseline limited the reach to long periods, and PLATO (2027) and Gaia DR4 (late 2026) will resolve the true occurrence.}
\end{frame}


% ════════════════════════════════════════════════════════════
\sectionimage{meadowsbarnes2018_habitability_factors}{Factors affecting planetary habitability. Meadows \& Barnes (2018)}
\section{Habitability as a Coupled Systems Property}
% ════════════════════════════════════════════════════════════

\begin{frame}
  \ipsherokey[0.58]{magnetosphere_anatomy_esa}{Anatomy of Earth's magnetosphere, one output of the geophysical engine. Credit: ESA, CC BY-SA 3.0 IGO.}{Habitability is a coupled stack, not a checklist: a stable, benign star and orbit; a geophysical engine that drives a dynamo and recycles volatiles; an atmosphere that keeps surface water liquid; a biosphere that feeds back on weathering and redox.}
\end{frame}

\begin{frame}
  \ipsherokey[0.65]{meadowsbarnes2018_habitability_factors}{Factors affecting habitability: stellar, planetary and orbital properties jointly shape the surface environment; blue observable, green model-interpreted, orange theoretical. Meadows \& Barnes (2018).}{Breaking any one of these couplings can end surface water stability.}
\end{frame}

\begin{frame}
  \ipsherokey{wordsworth2022_tidally_locked_planet}{Atmospheric circulation on a tidally locked rocky planet: substellar convection and equatorial super-rotation transport heat. Wordsworth \& Kreidberg (2022).}{Nightside cold traps risk atmospheric collapse.}
\end{frame}

\begin{frame}
  \ipsherokey[0.66]{kopparapu2014_hz_stellartype}{Habitable zone limits against stellar effective temperature. Kopparapu et al.\ (2014).}{Inner edge: the runaway greenhouse limit; outer edge: the maximum CO$_2$ greenhouse; stars brighten over their main-sequence lifetime, so orbiting in the zone today does not guarantee a habitable history.}
\end{frame}

\begin{frame}
  \ipsherokey[0.65]{lugerbarnes2015_runaway_duration}{Duration of the runaway greenhouse phase on a water-bearing planet against stellar mass. Luger \& Barnes (2015).}{Late M dwarfs spend up to a Gyr in the pre-main-sequence runaway: the planet starts \emph{inside} the runaway zone and loses its initial water inventory.}
\end{frame}


% ════════════════════════════════════════════════════════════
\sectionimage{lichtenberg2023_magma_ocean_differentiation}{Magma-ocean differentiation. Lichtenberg et al.\ (2023)}
\section{Magma Oceans, Water Delivery, Tectonic Thermostat}
% ════════════════════════════════════════════════════════════

\begin{frame}{Why "Delivery" Isn't the Right Frame}
  \begin{itemize}\setlength\itemsep{4pt}
    \item Earth, Venus, and Mars received similar initial accreted volatile budgets
    \item Magma ocean partitioning sets the core, mantle, and steam atmosphere reservoirs
    \item Evolutionary controls: early hydrodynamic escape, mantle outgassing, and stellar brightening
  \end{itemize}
  \vskip8pt
  \keyresult{Terrestrial volatile inventories diverge because differentiation and evolution matter far more than initial delivery.}
\end{frame}

\begin{frame}
  \ipsherokey{hamano2013_typeI_typeII}{Two types of terrestrial planet from magma ocean evolution. Credit: Hamano et al.\ (2013).}{Inside the critical flux, Type II planets lose their water to hydrodynamic escape before solidifying; Type I planets further out solidify quickly and retain oceans.}
\end{frame}

\begin{frame}
  \ipsherokey[0.61]{lichtenberg2023_magma_ocean_differentiation}{Magma ocean differentiation into a metal core, a silicate mantle and a steam atmosphere. Lichtenberg et al.\ (2023).}{Volatiles partition between liquid silicate and the steam atmosphere; the sequestered fraction depends on mantle redox, solidification time and depth, so identical initial budgets yield radically different atmospheric masses.}
\end{frame}

\begin{frame}
  \ipsherokey[0.66]{lammer2018_carbonate_silicate}{The carbonate-silicate cycle. Credit: Lammer et al.\ (2018).}{Volcanic CO$_2$ outgassing is balanced by temperature-dependent silicate weathering and subduction, a negative feedback that needs surface liquid water, active volcanism and tectonic recycling; its failure produces the Venus-like 92~bar.}
\end{frame}


% ════════════════════════════════════════════════════════════
\sectionimage{sandberg2018_drake_posterior}{Sandberg-Drexler-Ord Drake posterior (2018)}
\section{Drake Equation and Fermi Paradox}
% ════════════════════════════════════════════════════════════

\begin{frame}{The Drake Equation: A Framing Tool}
  \[
    N = R_\star \cdot f_p \cdot n_e \cdot f_l \cdot f_i \cdot f_c \cdot L
  \]
  $R_\star$ star-formation rate, $f_p$ fraction with planets, $n_e$ habitable planets per system, $f_l$ life, $f_i$ intelligence, $f_c$ communication, $L$ lifetime of a signalling civilisation.
  \begin{itemize}
    \item Astrophysical factors are now measured: $R_\star$, $f_p \sim 1$, $n_e \approx 0.1$--$0.6$ (Bryson+ 2021)
    \item Biological and social terms ($f_l, f_i, f_c, L$) span $4$--$10$ orders of magnitude
    \item Factors co-evolve rather than acting as independent pipeline stages
  \end{itemize}
  \vskip6pt
  \keyresult{The Drake equation is a tool for framing ignorance, not a numerical estimator.}
\end{frame}

\begin{frame}
  \ipsherokey[0.65]{sandberg2018_drake_posterior}{Monte Carlo posterior for the number $N$ of communicating civilisations, sampled over the full published parameter ranges. Sandberg, Drexler \& Ord (2018).}{The distribution is \emph{bimodal}, with substantial probability below $N = 1$ in the Milky Way: the apparent cosmic silence is unsurprising once the priors are honest.}
\end{frame}

\begin{frame}{Why Bimodal? Log-Sum of Broad Priors}
  Taking the logarithm turns the Drake product into a sum of broad priors:
  \[
    \log_{10} N = \log_{10}(R_\star f_p n_e) + \log_{10} f_l + \log_{10} f_i + \log_{10} f_c + \log_{10} L
  \]
  \begin{itemize}
    \item Unconstrained factors have log-uniform priors spanning many decades
    \item The log-sum of broad, skewed priors spreads $\log N$ over tens of decades; the long low tail of $f_l$ makes the result bimodal
    \item The posterior spans $\sim 30$ decades, placing high probability on $N < 1$
  \end{itemize}
\end{frame}

\begin{frame}
  \ipsherokey[0.62]{fermi_filters}{The Fermi paradox as a chain of filters, with a galactic crossing of $10^6$ to $10^7$~yr against a $10^{10}$~yr galaxy age. Course figure.}{The silence is an intuition check on subjective priors, not physical evidence: rare abiogenesis or intelligence, a Great Filter that keeps $L$ short, or transmissions below current detection limits.}
\end{frame}


% ════════════════════════════════════════════════════════════
% BREAK
% ════════════════════════════════════════════════════════════
\breakslide[End of Part 1: Solar System Context and Habitability\\Part 2: Habitable Zone Derivation, Astrobiology, Life Detection]


% ════════════════════════════════════════════════════════════
\sectionimage{kopparapu2014_hz_stellartype}{Habitable zone vs $T_\mathrm{eff}$. Kopparapu (2014)}
\section{Blackboard: Habitable Zone Boundaries}
% ════════════════════════════════════════════════════════════

\begin{frame}{Goal and strategy}
  \begin{columns}[c]
    \column{\recapfigcolnarrow}
    \ipsrecapfig[0.62\textheight]{board_sketch_l14}
    \column{\recaptextcolwide}
    \small
    \textbf{Goal:} locate the inner and outer edges of the habitable zone from one energy balance.
    \vskip4pt
    \begin{itemize}\setlength\itemsep{2pt}
      \item Inverse square in, Stefan-Boltzmann out: $T_\mathrm{eq}(d)$
      \item Express the boundaries as an effective flux $S_{\mathrm{eff}}$
      \item Insert the runaway limit (Lecture 9) and the maximum CO$_2$ greenhouse
    \end{itemize}
    \vskip4pt
    \keyresult{Star $L_\star$, planet at $d$ with Bond albedo $A_B$: the zone is the strip between $d_{\mathrm{in}}$ and $d_{\mathrm{out}}$.}
  \end{columns}
\end{frame}

\begin{frame}{Setup: Stellar Flux and Equilibrium Temperature}
  Star with bolometric luminosity $L_\star$, planet at distance $d$:
  \[
    F_\star(d) = \frac{L_\star}{4\pi d^2}
  \]
  Bond albedo $A_B$ (reflected fraction), rapidly rotating sphere, equilibrium:
  \[
    T_\mathrm{eq}(d) = \left[\frac{L_\star (1 - A_B)}{16\pi\sigma\epsilon\,d^2}\right]^{1/4}
  \]
  For Earth: $A_B = 0.3$, $\epsilon = 1$, $L_\star = \Lsun$, $d = 1$ AU $\to$ $T_\mathrm{eq} \approx 255$ K; the surface is 288 K from $\sim 33$ K greenhouse warming. \textit{To the board.}
\end{frame}

\begin{frame}{Result: Habitable Zone Boundaries}
  Effective flux $S_{\mathrm{eff}} = F_\star(d)/S_\oplus$ and the boundary distances:
  \[
    d = 1\ \mathrm{AU}\,\sqrt{\frac{L_\star/\Lsun}{S_{\mathrm{eff}}}}
  \]
  \begin{itemize}
    \item Inner edge: runaway greenhouse (Simpson-Nakajima, Lecture 9), $S_{\mathrm{in}} \approx 1.06$, $\approx 0.97$ AU for the Sun
    \item Outer edge: maximum $\mathrm{CO_2}$ greenhouse, $S_{\mathrm{out}} \approx 0.35$, $\sim 1.69$ AU for the Sun
    \item Main-sequence lifetime $t_{\mathrm{MS}} \approx 10\ \mathrm{Gyr}\,(M_\star/M_\odot)^{-2.5}$ sets how long the zone lasts
  \end{itemize}
\end{frame}

\begin{frame}{HZ Across Stellar Types}
  Habitable zone distance scales with stellar luminosity: $d \propto \sqrt{L_\star}$.
  \vskip6pt
  \begin{itemize}\setlength\itemsep{4pt}
    \item \textbf{G dwarf} ($\Lsun$): $d = 0.97$--$1.69$ AU; main-sequence lifetime $\sim 10$ Gyr
    \item \textbf{K dwarf} ($0.3\,\Lsun$): $d = 0.53$--$0.93$ AU; lifetime $24$ Gyr (orange-dwarf sweet spot)
    \item \textbf{M dwarf} ($5\times 10^{-4}\,\Lsun$): $d = 0.022$--$0.038$ AU; lifetime $10^{12}$ yr, but tidally locked with a long pre-MS runaway
  \end{itemize}
\end{frame}

\begin{frame}
  \ipsherokey[0.62]{cuntz_guinan_xray_lyman}{Lyman-$\alpha$ and X-ray fluxes at the Earth-equivalent distance against stellar type and age. Cuntz \& Guinan (2016).}{M dwarfs deliver over 100 to 500 times the G-dwarf flux; X-rays drop 500-fold over 5~Gyr but Lyman-$\alpha$ only 5 to 10-fold, so early K dwarfs offer long lifetimes without the M-dwarf XUV extremes.}
\end{frame}


% ════════════════════════════════════════════════════════════
\sectionimage{yang2013_tidally_locked_clouds}{Tidally locked planet clouds. Yang et al.\ (2013)}
\section{HZ Corrections: 3D Climate, Tidal Locking}
% ════════════════════════════════════════════════════════════

\begin{frame}
  \ipsherokey{yang2013_tidally_locked_clouds}{3D global climate model of a tidally locked planet with clouds. Yang et al.\ (2013).}{Persistent substellar convection lifts thick water clouds, the planetary albedo rises to $\sim 0.6$, and the inner edge tolerates $\sim 2\times$ the flux that 1D models predict.}
\end{frame}

\begin{frame}
  \ipsherokey[0.61]{shields2016_hz_tidally_locked}{Habitable-zone limits for M dwarfs with stellar flux increasing to the left: 1D Kopparapu boundaries (solid), the 3D inner edge (red dotted) and the tidal-locking radius (grey dashed). Shields et al.\ (2016).}{The 3D inner edge tolerates $1.7$--$2.2\times$ the 1D flux, and the tidal-locking radius encloses almost all M-dwarf habitable-zone planets.}
\end{frame}

\begin{frame}{Tidal-Locking Timescale}
  Tidal despinning scales steeply as $a^6$, so close-in habitable planets around low-mass stars synchronise quickly:
  \[
    \tau_\mathrm{lock} \sim \frac{\omega_0\,a^6\,I_p\,Q_p}{3\,G\,M_\star^2\,k_{2,p}\,R_p^5}
  \]
  \begin{itemize}
    \item Earth at 1 AU has $\tau_\mathrm{lock} \sim 10^{12}$ yr, far exceeding the age of the universe
    \item TRAPPIST-1 b ($a \approx 0.011$ AU) synchronises in $\lesssim 10^7$ yr
    \item With a system age of $\sim 8$ Gyr, all known M-dwarf HZ worlds are locked
  \end{itemize}
\end{frame}


% ════════════════════════════════════════════════════════════
\sectionimage{schwieterman2018_biosignature_classes}{Biosignature classes. Schwieterman et al.\ (2018)}
\section{Astrobiology and Biosignatures}
% ════════════════════════════════════════════════════════════

\begin{frame}{Origin-of-Life Hypotheses}
  \begin{itemize}\setlength\itemsep{4pt}
    \item \textbf{RNA world}: catalytic ribozymes carry the genetic code before proteins
    \item \textbf{Alkaline hydrothermal vents}: geochemical proton gradients drive early metabolism
    \item \textbf{Warm surface ponds}, where wet-dry cycles polymerise organic precursors, and \textbf{panspermia}, the transfer of viable organisms between bodies
  \end{itemize}
\end{frame}

\begin{frame}
  \ipsherokey{life_definition}{Working definition of life, with water and carbon as defaults and other solvents and chemistries not excluded. Course figure.}{Life (Lederberg 1965): a self-replicating, metabolising, evolving chemical system.}
\end{frame}

\begin{frame}
  \ipshero{earth_eons_timeline}{Geologic eons and eras of Earth across 4.54~Gyr: the Precambrian occupies almost 90\% of the timeline, with life emerging early in the Archean. Credit: Course-original figure; Gradstein et al.\ (2020).}
\end{frame}

\begin{frame}
  \ipsherokey[0.66]{schwieterman2018_biosignature_classes}{Three classes of biosignature. Schwieterman et al.\ (2018).}{\textbf{Gaseous} (O$_2$, O$_3$, CH$_4$, N$_2$O, organosulfur species), \textbf{surface} (the vegetation red edge near 700~nm) and \textbf{temporal} (seasonal cycles such as Earth's CO$_2$ Keeling oscillation): a credible detection needs several classes together.}
\end{frame}

\begin{frame}
  \ipsherokey[0.66]{schwieterman2018_biosignature_gases}{Mid-infrared bands of ten candidate biosignature gases. Schwieterman et al.\ (2018).}{O$_3$ at $9.6\,\mu$m, CH$_4$ at 3.3 and $7.7\,\mu$m, N$_2$O and the DMS and DMDS window in 3--15~$\mu$m set the wavelength needs of LIFE (mid-IR interferometer) and HWO (Habitable Worlds Observatory).}
\end{frame}

\begin{frame}{Disequilibrium is the Diagnostic}
  \begin{itemize}\setlength\itemsep{4pt}
    \item Individual gases in isolation almost always possess plausible abiotic origins
    \item Earth pair: $\mathrm{O_2}$ (21\%) and $\mathrm{CH_4}$ (1.8 ppm) react on decade timescales
    \item Simultaneous persistence requires massive, continuous biological resupply fluxes
  \end{itemize}
  \vskip8pt
  \keyresult{Chemically incompatible gases coexisting in an atmosphere are the most convincing remote signature of active biology.}
\end{frame}

\begin{frame}
  \ipsherokey{madhusudhan2023_k218b_spectrum}{JWST transmission spectrum of K2-18~b combining NIRISS and NIRSpec observations. Credit: Madhusudhan et al.\ (2023).}{Absorption bands reveal atmospheric CH$_4$ and CO$_2$ alongside a tentative DMS feature near $3.4\,\mu$m.}
\end{frame}

\begin{frame}{Catling 2018: A Bayesian Framework}
  \begin{columns}[T]
    \column{0.5\textwidth}
    \includegraphics[width=\linewidth, height=0.52\textheight, keepaspectratio]{catling2018_assessment_framework}
    \column{0.5\textwidth}
    \includegraphics[width=\linewidth, height=0.52\textheight, keepaspectratio]{catling2018_bayesian_framework}
  \end{columns}
  \vskip4pt
  \keyresult{Detection is a Bayesian chain of inference over stellar context, planetary properties and abiotic mechanisms, not a single yes/no spectral peak.}
  \sourcenote{Catling et al.\ (2018)}
\end{frame}


% ════════════════════════════════════════════════════════════
\sectionimage{lincowski2021_phosphine_so2}{Venus phosphine controversy. Lincowski et al.\ (2021)}
\section{False Positives and Solar System Targets}
% ════════════════════════════════════════════════════════════

\begin{frame}{False Positives: Every Gas Has an Abiotic Pathway}
  \begin{itemize}\setlength\itemsep{4pt}
    \item Abiotic $\mathrm{O_2}$: $\mathrm{H_2O}$ photolysis followed by hydrogen escape, or $\mathrm{CO_2}$ photolysis
    \item Abiotic $\mathrm{CH_4}$: hydrothermal serpentinisation or volcanic outgassing at low oxygen fugacity
    \item Abiotic $\mathrm{N_2O}$ from lightning and photochemistry; DMS has no significant abiotic Earth source, but one could exist elsewhere
  \end{itemize}
  \vskip8pt
  \keyresult{Verification is an inverse problem: confirmation requires ruling out every plausible geochemical and photochemical false positive.}
\end{frame}

\begin{frame}
  \ipsherokey{madhusudhan_k218b_dms_post}{Retrieved mixing ratio posteriors for CH$_4$, CO$_2$ and DMS in K2-18~b under zero, one and two floating instrument offsets. Credit: Madhusudhan et al.\ (2023).}{CH$_4$ and CO$_2$ hold under every offset choice, while the DMS significance collapses from $2.4\sigma$ to below $1\sigma$.}
\end{frame}

\begin{frame}{Solar System Targets for Life Detection}
  \begin{itemize}\setlength\itemsep{4pt}
    \item Mars: ancient habitable lacustrine environments; the Mars Sample Return path
    \item Europa and Enceladus: confirmed subsurface oceans; Enceladus's plume carries $\mathrm{H_2}$, organics and phosphates from hydrothermal activity, hypothesised for Europa
    \item Titan and Venus: complex hydrocarbon hydrology and contested cloud chemistry
  \end{itemize}
  \vskip8pt
  \keyresult{Five prime environments with liquid water reservoirs or rich organic chemistry.}
\end{frame}

\begin{frame}
  \ipshero[0.76]{extremophile_envelope}{Extremophiles widen the definition of habitable: Earth life spans $-20$ to over $120\,^\circ$C, and the icy moons hold liquid water in contact with rock outside the classical habitable zone. Course figure.}
\end{frame}

\begin{frame}
  \ipsherokey{lincowski2021_phosphine_so2}{The Venus phosphine controversy: the claimed PH$_3$ line against the mesospheric SO$_2$ fit. Lincowski et al.\ (2021).}{Greaves et al.\ (2021) claimed PH$_3$ in the Venus clouds, but the line is fit by SO$_2$; K2-18~b teaches the same lesson (CH$_4$ at $5\sigma$ and CO$_2$ at $3\sigma$ stand, DMS fell from $2.4\sigma$ to $\sim 1\sigma$); DAVINCI tests Venus in situ, EnVision and VERITAS from orbit.}
\end{frame}


% ════════════════════════════════════════════════════════════
\sectionimage{greene2023_trappist1b_eclipse}{JWST/MIRI eclipse of TRAPPIST-1 b. Greene et al.\ (2023)}
\section{Exoplanet Life Detection}
% ════════════════════════════════════════════════════════════

\begin{frame}
  \ipsherokey{greene2023_trappist1b_eclipse}{JWST MIRI $15\,\mu$m secondary eclipse of TRAPPIST-1~b. Credit: Greene et al.\ (2023).}{The dayside brightness temperature of $503$~K matches the $508$~K bare-rock prediction and rules out a thick atmosphere.}
\end{frame}


% ════════════════════════════════════════════════════════════
\section*{Course Wrap-up}
% ════════════════════════════════════════════════════════════

\begin{frame}
  \ipshero{five_lessons}{The five biggest lessons: formation, interiors as heat engines, dynamic atmospheres, habitability as a coupled property, the solar system as reference case. Course figure.}
\end{frame}

\begin{frame}
  \ipshero{five_questions}{The five biggest open questions, each with what is at stake or the test that will decide it. Course figure.}
\end{frame}

\begin{frame}{Summary}
  \begin{enumerate}\setlength\itemsep{6pt}
    \item Planet formation and differentiation are universal physical processes; interiors are heat engines that drive tectonics, dynamos and outgassing over Gyr
    \item Atmospheres evolve through escape, outgassing and photochemistry, and habitability is a coupled trajectory, not a static orbital condition
    \item The solar system is one detailed reference in a diverse exoplanet population
  \end{enumerate}
  \vskip8pt
  \textbf{Next:} The final exam.
\end{frame}

% ── Final slide ─────────────────────────────────────────────
\begin{frame}[plain,noframenumbering]
  \begin{tikzpicture}[remember picture, overlay]
    \ipscontain{title_background}
    \fill[ipsPrimary, opacity=0.70] (current page.north west) rectangle (current page.south east);
  \end{tikzpicture}
  \vfill
  \begin{center}
    {\color{white}\Huge \textbf{Thank You}}\\[14pt]
    {\color{white}\large Planetary Systems is complete.}\\[10pt]
    {\color{white}\normalsize Lecture notes: \texttt{formingworlds.github.io/IntroductionToPlanetaryScience}}
  \end{center}
  \vfill
\end{frame}

\end{document}
